%% kapa: exam=1 round=35 year=2024 period=1 form=A subject=civil_law subject_ko=민법 seq=33 number=33
%% source: Q-Net 감정평가사 기출문제 (cst00309) artlSeq=5213229 file=raw/1008/5213229/1교시 문제지 원본.pdf
%% conversion: kordoc 4.24.51 -> tools/gen1.py (auto); review=proofread-sonnet; flags=none
\begin{question}{33}
지상권에 관한 설명으로 옳지 않은 것은? (다툼이 있으면 판례에 따름)
\begin{choices}
  \item 저당물의 담보가치를 유지하기 위해 설정된 지상권은 피담보채권이 소멸하면 함께 소멸한다.
  \item 기존 건물의 사용을 목적으로 설정된 지상권은 그 존속기간을 30년 미만으로 정할 수 있다.
  \item 수목의 소유를 목적으로 하는 지상권이 존속기간의 만료로 소멸한 경우, 특약이 없는 한 지상권자가 존속기간 중 심은 수목의 소유권은 지상권설정자에게 귀속된다.
  \item 양도가 금지된 지상권의 양수인은 양수한 지상권으로 지상권설정자에게 대항할 수 있다.
  \item 토지양수인이 지상권자의 지료 지급이 2년 이상 연체되었음을 이유로 지상권소멸청구를 하는 경우, 종전 토지소유자에 대한 연체기간의 합산을 주장할 수 없다.
\end{choices}
\end{question}
\answer{3}
